\subsection*{Mathematical purpose}
A journey is an ordered finite list of directed traversals of actual edges on a fixed graph. The only shortening move deletes two adjacent traversals of the same edge in opposite directions. The graph, start, and finish remain fixed. The empty journey means staying at the initial vertex.

\textbf{Unique reduction.} Every journey ends in the same reduced edge path, regardless of deletion order. Each deletion shortens by two, so the process terminates. A running stack of oriented edges gives an invariant normal form; the proof is on page 6. ``Same endpoints'' is enough to identify an edge only on these simple printed graphs. Parallel edges would need separate road identities.

\textbf{Tree theorem.} Between two vertices of a tree there is a unique simple path, and every journey between them reduces to it. In particular, every closed tree journey reduces to the empty journey. A reduced path on a four-edge ring, by contrast, keeps going in a single direction and may survive for one or more complete circuits.

\textbf{Order can survive.} Let a be H$\to$A$\to$B$\to$H and b be H$\to$C$\to$D$\to$H on the two-ring map. The journeys ab and ba have distinct reduced paths. The commutator ab a$^{-1}$ b$^{-1}$ has zero net count around each ring but has no adjacent inverse edges, so it does not disappear. Net traversal counts lose essential order information.

\textbf{Topological scope.} These elementary edge moves describe endpoint-fixed homotopy of graph edge paths, with immediate inverse-pair insertions allowed when comparing equivalence classes. The student shortening activity uses deletions. The open regions inside rings are not filled faces: paths cannot slide across them, cancel separated inverse steps, trade subjourney order, or move their basepoint. Paths may revisit vertices and overlap themselves.

\subsection*{Entry points and readiness}
\textbf{K-1:} play and replay short journeys with a pawn; an adult records ordered steps. \textbf{Grades 2-3:} construct long tree walks, classify ring results, and contrast order on two rings. \textbf{Grades 4-5:} prove unique reduction and find the zero-net-count obstruction. Counting to 12 is useful, not fluent reading or algebra.

\textbf{Before-use gate:} the exact large route tiles, adult-recording method, mats, and pawn workflow have not been physically tried with the intended group. A K-1 child must be able to replay an ordered route and remove an immediate reversal without changing the remaining order. Stay with the concrete demonstration if that is not yet reliable.

\subsection*{Reading the keys}
In this key a compact route such as ABDE means A$\to$B$\to$D$\to$E. Count edges, not letters: a seven-step route has eight displayed vertices. The starting letter remains when every step is removed.
